% methods.tex -- Section III-A (pipeline under audit) and III-B (audit checks). III-C to III-E
% are in methods_def.tex, input on the last line. Labels: sec:methods, sec:pipeline, sec:audit,
% eq:fmedial.
% Sources (facts_verified.md IDs unless stated; M05 is WRONG and not used):
%   A: phi 0.5, d 45 mm, 1.4 BW (A1-8, ISRR8); BW = m g, Mx / (BW H) (M25); median height
%      1.712 m (M23); Rajagopal 37 DOF / 80 muscles (M02); SO settings (M04); reserves (M06);
%      0.3 s buffer (M08); pool and test split (M11, M13, ISRR7); filters (M17, M18);
%      inputs (M26-M30); architecture and loss (M31, M32, M34). Forward pass: our own timing,
%      2026-11-10 (not printed since the trim), fpopt checkpoint (235 inputs, 64,516 parameters), PyTorch 2.11, one thread,
%      Intel Core Ultra 9 285K, one query incl. standardisation, mean of 5,000 calls after 500
%      warm-up calls: 0.068-0.071 ms in three repeats (M43's 0.08 ms is [4]'s own claim, not used).
%      OpenSim version at labelling time is unverified, so no version number is printed.
%      The audited checkpoints follow the recipe [4] describes (235 inputs, no pelvis translations,
%      F_medial loss term). Which checkpoint the optimizer runs of [4] used is not confirmed
%      (facts ISRR4: possibly mlp-grffree-base, 198 inputs with pelvis translations, plain MSE);
%      TODO before submission: confirm with the first author of [4] which checkpoint produced
%      its figures. Access: the labels, label/training/inference code, two checkpoints and the
%      reference motion come from the authors of [4] (true because we are among them; check
%      wording against ARSO's double-blind rule, see \anonselftrue in main.tex). The trajectory
%      converter (pipeline/label_trajectory_csv.py, added 2026-09-23) was written for this study, not
%      by [4]; results.tex calls it 'the converter we used'.
%   B: empty forces_file, OpenSim 4.5.2 behaviour, vasti check, relabel of 9,529 folders
%      (gc/pilot_AB_sojr/README.md "OpenSim label sources"; qc_final.md fix:label-qc,
%      fix:train-plumbing); reference motion has jcf_output_sojr (lead, 2026-11-10);
%      sign checks (labels_qc/mx_sign/README.md); QC rules and thresholds (qc_final.md
%      build/fix:so-diagnosis); GC alignment, segments, regressions, mirroring, peak windows,
%      Kaneda reference and DM timing (gc/pilot_AB_sojr/README.md).
\section{Methods}\label{sec:methods}

\subsection{Pipeline Under Audit}\label{sec:pipeline}
\textit{Proxy.} Following \Chong{}~\citeChong{}, the pain-risk proxy is the
medial compartment load of a frontal-plane knee model~\cite{schipplein1991},
\begin{equation}\label{eq:fmedial}
  \Fmed = \phi\,\Fjcf + \Madd/d ,
\end{equation}
with $\Fjcf$ the compressive tibiofemoral force, $\Madd$ the adduction moment (positive when
loading the medial compartment), $\phi$ the medial share of $\Fjcf$ at $\Madd = 0$, and $d$
the distance between the compartment centres. Both loads come from an OpenSim JR analysis,
reported on the tibia in the tibial frame: $\Fjcf = |F_y|$ and $\Madd = s\,M_x$,
where the sign $s$ is audited in D2. With forces in BW and $M_x$ in BW$\cdot H$ (height $H$),
$\Fmed = \phi|F_y| + s M_x H/d$ in BW. Nominally $\phi = 0.5$,
$d = 0.045$~m and the limit is 1.4~BW~\citeChong{}; at the median pool height of
1.712~m, $H/d \approx 38$ amplifies moment errors.

\textit{Labels.} From AddBiomechanics~\cite{werling2024} the pipeline takes the kinematics, GRF
and centres of pressure (COP) and the scaled Rajagopal model~\cite{rajagopal2016} (37 degrees of
freedom, 80 muscles). SO minimizes summed squared activations (coordinates filtered at 6~Hz, at
most 200 iterations, reserve actuators on all unconstrained coordinates). The training pool holds
7,991 walking trials of 82 people from four labs: Carter~\cite{carter2024}, Moore~\cite{moore2015}, Lencioni~\cite{lencioni2019} and
Falisse~\cite{falisse2017}. The test split has one person from each of six labs (676 walking
trials), adding Fregly (Grand Challenge subject SC) and Han~\cite{han2023}, labs that supply no
training data. Trials are dropped if the resultant knee force exceeds 6~BW or $|M_x|$ exceeds
0.15~BW$\cdot H$, and training trials also if the waveform is flat.

\textit{Surrogate.} In the recipe that \Chong{} describe, a multilayer perceptron maps
$13 + 20\times11 + 2 = 235$ inputs to the
right-knee $(F_x, F_y, F_z, M_x)$: 13 joint angles at query frame $k$; for frames $k-10$ to
$k+9$, the whole-body centre of mass and a point on each calcaneus relative to the pelvis in
ground axes, and right and left contact flags (vertical GRF above 0.05~BW); and mass and height.
It takes no GRF and no pelvis translation. The foot point is the calcaneus origin or the
optimizer's foot point, $(0.13,-0.01,0)$~m in calcaneus axes. Three hidden layers of 128 units
are trained with Adam on a squared-error loss that includes 0.5 times the error of $\Fmed$ from
\eqref{eq:fmedial}. We audited two checkpoints trained with this recipe, one per foot point (`legacy').
The stored labels, the label, training and inference code, these checkpoints and the optimizer's
reference motion were made available by the authors of~\citeChong{}; the optimizer was not.

\subsection{Audit Checks}\label{sec:audit}
\textit{D1, label provenance.} The stored labels came from one OpenSim AnalyzeTool that ran SO
and JR with an empty forces file. JR then computes the muscle forces from the model states
instead of applying the SO solution. We compared the stored knee force with the SO force of the
vasti alone, and reran JR with the SO forces for all 9,529 labelled trials and for the
optimizer's reference motion, keeping the SO solution and time grid of the stored labels.

\textit{D2, moment sign.} JR reports the load on the tibia at the knee centre in the tibial
frame, whose $z$ axis points to the subject's right on both legs. We tested which $s$ makes
$\Madd$ medial-positive in three ways: contact algebra (a compressive condyle force $f$ at
$z_c$ gives $M_x = z_c f$), the $z$ positions of the medial and lateral knee-muscle points in one
scaled model per lab, and a fit $M_x = aF_{\mathrm{med}} + bF_{\mathrm{lat}} + c$ on Kaneda et al.'s
released SO+JR results~\cite{kaneda2023}, which include both condyle forces.

\textit{D3, SO label quality.} Rules on the SO and JR outputs (Table~\ref{tab:qc}) target three
defects. \emph{Fixed COP}: the right foot's COP moves slower than 0.2~m/s in gait stance, so its
GRF acts at a point fixed in the lab. \emph{Unconverged SO}: at sampled stance frames we re-solved
SO's convex problem from the stored moment arms and net moments; a stored cost above 1.1 times
the optimum marks a non-optimal frame. Frame-level proxies are saturation, reserves,
co-activation and activation jumps. \emph{Dynamic inconsistency}: a pelvis residual moment above
5\% of peak GRF times pelvis height. We also classified the mechanism of 40 hand-picked SO+JR
peaks above 4~BW.

\textit{D4, input construction.} We compared the time grids on which the legacy feature builder
places position windows, joint angles, flags and labels, and the time spanned by a window.

\textit{In vivo comparison.} The Knee Grand Challenge~\cite{fregly2012} provides implant forces
of subjects JW and DM (right knee) and SC (left knee), none used in training.
We matched 138 AddBiomechanics segments to Grand Challenge trials by cross-correlating vertical
GRF ($r \ge 0.98$, amplitude ratio 0.9--1.1), took the supplied medial force (JW: sum of the
medial load cells), and scored instrumented-limb stance frames inside the JR window at 100~Hz.
The surrogates received aligned 100~Hz inputs. For
SC, JR used the left knee and the surrogate inputs were mirrored. Peaks were scored in stance
windows at 5--45\% and 55--95\%. We also scored Kaneda et al.'s released
SO+JR~\cite{kaneda2023} on SC. DM's implant timing is
suspect (heel strike to measured first peak takes 95~ms longer than in gait-matched JW and SC
stances), so DM results are conditional.

% methods_def.tex -- Section III-C (sensitivity), III-D (corrected baseline), III-E (wearable).
% Input on the last line of methods.tex. Labels: sec:sens, sec:corrected, sec:wearable.
% Sources: S1 grid and limit (A1-1, A1-3); 1.3 cm pelvis lift, joints fixed (M36, PS6, PS7);
%   foot points (M29); S3 deployment form 20 frames at 100 Hz (A3-34); reference motion 6.45 s at
%   100 Hz (GT6); its SO+JR labels are unusable (label_mask keeps 0 right-knee frames; sep23/README).
%   Corrected baseline = user decision (a)/(b), 2026-11-10. Training set = train_mlp *_sojrqc input
%   sets: keep_r of <trial>/jcf_output_sojr/label_mask.npz (fixed right-foot COP trials, then the six
%   SO/dynamics frame rules; the 5 BW flag is stored, NOT applied), carried to the 100 Hz grid with
%   LABEL_MASK_TOL = 0.05, physical M_x sign (MX_SIGN_PHYSICAL). 25.4 % = kept_stance_frac_r of the
%   training walking pool in results/thesis/labels_qc/label_masks/by_lab.csv (adding the 5 BW flag
%   would remove another 2.5 % of the kept stance time, fy5_after_so_frac_of_kept_stance_r). The
%   legacy trial filters (6 BW, 0.15 BW*H) would drop 86 % of the kept training frames
%   (label_masks/dry_build_fpopt/trial_filters.csv), so the corrected runs switch them off
%   (--max-peak 1e9 --max-moment 1e9); keep in sync with the training commands.
%   Lencioni = the lab that earlier drafts called 'Tiziana' (data folder name). S2 pelvis lift: the numbers are [4]'s own, cited as such in results.tex; we did not repeat the
%   test (no checkpoint was lift-tested). S2 pelvis-angle conversion and S4: results/thesis/sep23
%   (converted pelvis angles are the primary inputs). S3: legacy in vivo, audit_gc_legacy_native vs
%   audit_gc. Seed 2 best epoch 8 (1-based), val 0.363; seeds 1/3 0.300/0.305 (sep23/README.md).
%   Simulated suit: two held-out trials, noise 0/1/2 deg and 1 deg + 0.5 deg/s drift (pn3_sim/README).
%   Wearable: A3-01..A3-03, A3-05, A3-22..A3-25, A3-29..A3-34, A3-39, A3-41, A3-07, A3-08;
%   Wang et al. data: results/thesis/twente/EXPLORE_NOTES.md sections 1-6 (9 people, 54 treadmill
%   trials at 0.25-1.50 m/s, 8 sensors, optical IK, treadmill force, separate static pose, no
%   knee-load labels).
\subsection{Sensitivity Protocol}\label{sec:sens}
\textit{S1, constants.} We recompute $\Fmed$ for $\phi \in \{0.4, 0.5, 0.6\}$ and
$d \in \{40, 45, 50\}$~mm, with the limit at 1.4~BW, and report the mean absolute error (MAE) and
the recall and precision of limit crossings. \textit{S2, inputs.} We swap the foot points and, on
the reference motion, compare two pelvis-angle conventions. \textit{S3, alignment.} Legacy models
are scored in vivo on training-style and on aligned 100~Hz inputs, as the optimizer and IMU
driver supply. \textit{S4, training run.} Models differing only in initialization seed are
scored on the test split, in vivo, and on the optimizer's reference motion~\citeChong{},
which lacks usable labels.

\subsection{Corrected Baseline}\label{sec:corrected}
The corrected baseline keeps the legacy architecture and recipe but (i)~uses SO+JR labels,
(ii)~uses the physical sign in the $\Fmed$ loss and stores it in the checkpoint, (iii)~builds
inputs and labels on the JR time grid resampled to 100~Hz, so a window spans 200~ms as in
deployment, and (iv)~trains only on frames that pass the quality rules (Table~\ref{tab:qc}),
which replace the legacy trial filters; a 100~Hz frame passes when at most 5\% of its
interpolated label comes from flagged frames. Seeds 1--3 use the optimizer's foot point and one
run (seed~1) the calcaneus origin, all on one person split. Seed~2 reached its best validation
loss at epoch~8 (0.363, against 0.300 and 0.305); we keep it. Legacy checkpoints keep their
formula, so results computed with them stay reproducible.

\subsection{Wearable Path}\label{sec:wearable}
A Perception Neuron~3 suit streams fused orientations of 8 sensors (pelvis, torso, thighs,
shanks, feet) at 60~Hz, but no pelvis position. After a 1-s standing calibration, an unscented
Kalman filter (UKF)~\cite{julier2004unscented} with $2\times38$ cubature
points~\cite{arasaratnam2009cubature} estimates 19 coordinates of a height-scaled Rajagopal model
and their velocities from orientation residuals ($\sigma = 2^\circ$), and legged
odometry~\cite{bloesch2012state} holds each planted foot ($\sigma = 1.5$~cm). Hysteresis detectors
on foot angular speed~\cite{skog2010zero} set the contact flags; as they switch on at foot flat,
about 90~ms after heel strike, each onset is back-dated by up to 0.2~s to the speed minimum at
initial contact. A 20-frame, 100~Hz window feeds the surrogate (90~ms lag); results go over Redis
to a visualizer.

We test the path on a simulated suit (orientations from two held-out trials with 0, 1 and
2$^\circ$ noise, and 1$^\circ$ with 0.5$^\circ$/s heading drift) and on the Xsens recordings of
Wang et al.~\cite{wang2023twente}: 9 people, 54 treadmill trials at 0.25--1.50~m/s, 8 sensors at
the suit's locations, optical motion capture, treadmill force and a separate static calibration
pose. Lacking knee-load labels, we score the IMU path against optical inverse kinematics (IK) and
against the same surrogate run on optical kinematics.

